% The argument on one page. Build:  latexmk -pdf one-page.tex
\PassOptionsToPackage{obeyspaces}{url}
\documentclass[10pt,a4paper]{article}
\usepackage[T1]{fontenc}
\usepackage[utf8]{inputenc}
\DeclareUnicodeCharacter{2082}{\textsubscript{2}}
\usepackage{lmodern}
\usepackage[final]{microtype}
\usepackage[a4paper,margin=0.62in,top=0.5in,bottom=0.5in]{geometry}
\usepackage{amsmath,amssymb}
\usepackage{xcolor}
\usepackage[colorlinks=true,urlcolor=blue!50!black]{hyperref}
\DeclareUrlCommand\lean{\urlstyle{tt}\def\UrlFont{\ttfamily\footnotesize}}
\newcommand{\Z}{\mathbb{Z}}\newcommand{\Q}{\mathbb{Q}}\newcommand{\F}{\mathbb{F}}
\newcommand{\fP}{\mathfrak{P}}
\newcommand{\re}{\operatorname{Re}}\newcommand{\Frob}{\operatorname{Frob}}
\newcommand{\ph}{\operatorname{ph}}\newcommand{\len}{\operatorname{len}}
\newcommand{\GL}{\mathrm{GL}}\newcommand{\U}{\mathrm{U}}\newcommand{\Sha}{\mathrm{Sha}}
\newcommand{\ASK}{\mathrm{ASK}}\newcommand{\MAIN}{\mathrm{MAIN}}\newcommand{\Br}{\mathrm{Br}}
\newcommand{\eps}{\varepsilon}\newcommand{\abs}[1]{\lvert #1\rvert}
\newcommand{\step}[1]{\par\smallskip\noindent\textbf{#1}}
% Version and date. The measured numbers are in numbers.tex, which tools/numbers.py writes from
% the results files. A few numbers are typed here because no results file holds them; each
% says where it comes from. numbers.tex also sets \ifsecondrun: false while the results are
% those of the first run only.
\newcommand{\paperdate}{23 November 2026}
\newcommand{\paperversion}{1.5}
% Generated by tools/numbers.py from verify/results/ and refs/rows-counts.json.
% Do not edit by hand.
\newif\ifsecondrun
\secondrunfalse
\newcommand{\filecount}{175{,}187}% scan.txt: lean files
\newcommand{\linecount}{17{,}374{,}215}% scan.txt
\newcommand{\linemillions}{17}% computed
\newcommand{\archivebytes}{239{,}275{,}554}% scan.txt
\newcommand{\archivesha}{9631b4f020199e580ba86178cb56c6e9b09f690c2369a80dc4334aa93f49c7fc}% scan.txt
\newcommand{\archiveshaA}{9631b4f020199e580ba86178cb56c6e9}% first half, for display
\newcommand{\archiveshaB}{b09f690c2369a80dc4334aa93f49c7fc}% second half, for display
\newcommand{\otherfiles}{3}% scan.txt
\newcommand{\mainlines}{12}% scan.txt
\newcommand{\scanmacros}{87}% scan.txt
\newcommand{\scanelabs}{4}% scan.txt
\newcommand{\scanruncmd}{5}% scan.txt
\newcommand{\scanelabfiles}{3}% scan.txt: files listed under elab
\newcommand{\scanruncmdfiles}{5}% scan.txt: files listed under run_cmd
\newcommand{\scanpartial}{2}% scan.txt
\newcommand{\scandepth}{203}% scan.txt
\newcommand{\scanimports}{909{,}822}% scan.txt
\newcommand{\scanallmathlib}{136{,}003}% scan.txt
\newcommand{\buildhours}{13.4}% build-summary.txt: END minus START
\newcommand{\buildjobs}{183{,}893}% build-summary.txt
\newcommand{\buildthreads}{32}% build-summary.txt
\newcommand{\buildcpu}{AMD EPYC 9575F 64-Core Processor}% build-summary.txt
\newcommand{\leanversion}{4.33.1}% build-summary.txt
\newcommand{\mathlibrev}{0df444a360eaa60ab8c11dca51a86af692955474}% build-summary.txt
\newcommand{\conesize}{1{,}478{,}219}% cut-result.txt
\newcommand{\conefive}{1{,}478{,}243}% cut-result.txt
\newcommand{\conefiveextra}{24}% difference of the two
\newcommand{\conearchive}{1{,}265{,}256}% the lists: names of the archive in the cone
\newcommand{\conerest}{212{,}963}% CONE-SIZE minus the archive's names in the cone
\newcommand{\projnames}{1{,}905{,}922}% the lists, by name
\newcommand{\projnamesout}{640{,}666}% the lists, by name
\newcommand{\projtheorems}{1{,}548{,}656}% the lists, by name
\newcommand{\coneprojtheorems}{1{,}075{,}015}% the lists, by name
\newcommand{\unusedtheorems}{473{,}641}% difference
\newcommand{\repeatedtheorems}{39{,}597}% theorem declarations minus theorem names
\newcommand{\theoremshare}{69\%}% computed
\newcommand{\projdefs}{331{,}347}% the lists, by name
\newcommand{\coneprojdefs}{178{,}298}% the lists, by name
\newcommand{\defshare}{54\%}% computed
\newcommand{\projunsaferec}{621}% the lists: definitions whose name ends in _unsafe_rec
\newcommand{\projmodules}{175{,}014}% cut-result.txt
\newcommand{\coneprojmodules}{147{,}489}% cut-result.txt
\newcommand{\moduleshare}{84\%}% computed
\newcommand{\unusedmodules}{27{,}698}% Lean files that declare nothing of the cone
\newcommand{\emptymodules}{173}% Lean files that declare nothing
\newcommand{\hubtheorems}{198}% cut-result.txt
\newcommand{\projopaque}{2}% the lists: opaque constants of the archive
\newcommand{\coneprojopaque}{0}% the lists: of them in the cone
\newcommand{\projinductives}{8{,}509}% the lists: inductive types of the archive
\newcommand{\replaymodules}{186{,}200}% replay-summary.txt
\newcommand{\replayconstants}{2{,}701{,}132}% replay-summary.txt
\newcommand{\replaydecls}{2{,}695{,}292}% replay-summary.txt
\newcommand{\replayskipped}{5{,}840}% replay-summary.txt
\newcommand{\replaytheorems}{2{,}084{,}256}% replay-summary.txt
\newcommand{\replaydefs}{556{,}318}% replay-summary.txt
\newcommand{\replayopaque}{3{,}882}% replay-summary.txt
\newcommand{\replayaxioms}{7}% replay-summary.txt
\newcommand{\replayquot}{4}% computed: the constants replayed less those of the other kinds; checked to be 4
\newcommand{\replayinductives}{15{,}769}% replay-summary.txt
\newcommand{\replayblocks}{15{,}734}% replay-summary.txt
\newcommand{\replayctors}{19{,}159}% replay-summary.txt
\newcommand{\replayrecs}{15{,}897}% replay-summary.txt
\newcommand{\replayminutes}{57}% replay-summary.txt: seconds / 60
\newcommand{\replayrepeated}{40{,}595}% replay-summary.txt
\newcommand{\replayskippedrec}{5{,}008}% replay-skipped.txt
\newcommand{\replaymillions}{2.7}% computed
\newcommand{\replayprogramlines}{582}% verify/replay/ReplayMain.lean: lines, comments included
\newcommand{\firstprogramlines}{294}% results/replay-program-first-form.lean: lines
\newcommand{\hostileprogramsha}{b5e3596dde657bc8}% hostile.txt: first sixteen digits
\newcommand{\hostilemodules}{29}% hostile.txt
\newcommand{\hostilebad}{25}% hostile.txt: modules a replay must reject
\newcommand{\hostilecontrols}{4}% hostile.txt
\newcommand{\hostilecrafted}{11}% hostile/: modules with no source file
\newcommand{\hostilefreshrejects}{9}% hostile.txt
\newcommand{\hostiledefaultrejects}{7}% hostile.txt
\newcommand{\independentfreshrejects}{9}% independent-hostile.txt
\newcommand{\independentdefaultrejects}{7}% independent-hostile.txt
\newcommand{\independentfreshminutes}{77}% independent-fresh.txt
\newcommand{\stockfreshrzminutes}{11}% stock.txt
\newcommand{\stockfreshrzcpuminutes}{11}% stock.txt: processor time
\newcommand{\stockreplayrzdecls}{430{,}239}% stock.txt
\newcommand{\stockreplayrzminutes}{1}% stock.txt
\newcommand{\stockreplayrzcpuminutes}{8}% stock.txt: processor time
\newcommand{\stockfreshmathlibminutes}{36}% stock.txt
\newcommand{\stockfreshmathlibcpuminutes}{36}% stock.txt: processor time
\newcommand{\stockreplaymathlibdecls}{766{,}714}% stock.txt
\newcommand{\stockreplaymathlibminutes}{5}% stock.txt
\newcommand{\stockreplaymathlibcpuminutes}{31}% stock.txt: processor time
\newcommand{\stockexportedmodules}{10{,}498}% replay-Mathlib/replay-summary.txt
\newcommand{\stockexported}{644{,}284}% replay-Mathlib/replay-summary.txt
\newcommand{\stockpartialunmarked}{531}% replay-Mathlib/replay-summary.txt
\newcommand{\toolchainfiles}{17{,}493}% toolchain.txt
\newcommand{\literaturerows}{186}% refs/rows-counts.json
\newcommand{\rowsdistinct}{184}% refs/rows-counts.json
\newcommand{\rowsrepeated}{2}% computed
\newcommand{\rowspublished}{179}% computed
\newcommand{\rowsgeneral}{90}% refs/rows-counts.json
\newcommand{\rowsspecial}{86}% refs/rows-counts.json
\newcommand{\rowsother}{3}% refs/rows-counts.json
\newcommand{\rowsinternal}{3}% refs/rows-counts.json

% Billed by the provider for the machine of the first run, from its start to its end; not a
% measurement of the archive.
\newcommand{\firstruncost}{26.81}
% The first replay was started with LEAN_NUM_THREADS=28; this is our record of the command,
% and no results file of that run holds it.
\newcommand{\firstreplaythreads}{28}
% The cost of the second machine, \secondruncost, is computed in tools/numbers.py from the
% provider's hourly price and the machine's start, which are typed there, and the last time in
% file-times.txt.
% The memory limit of the first machine's container, in GB, as the provider listed it; no
% results file of that run holds it.
\newcommand{\firstmemorylimit}{128}

% The outcome of the kernel replay, in words (used in the section on verification).
\newcommand{\replaystatus}{%
The kernel accepted every declaration sent to it: \replaytheorems{} theorems, \replaydefs{}
definitions, \replayopaque{} opaque constants, \replayaxioms{} axioms, the quotient type with
its \replayquot{} constants, and \replayinductives{} inductive types in \replayblocks{} blocks.
\ifsecondrun
The comparison found every one of the \replaydecls{} constants of the files, among them
\replayctors{} constructors and \replayrecs{} recursors, equal to the kernel's, and nothing
else in the kernel's environment. There was no failure. The run took \replayminutes{} minutes
on \replaythreads{} threads.
\else
The \replayctors{} constructors and \replayrecs{} recursors in the files agree with the ones
the kernel generated. There was no failure. The run took \replayminutes{} minutes on
\firstreplaythreads{} threads.
\fi
The \replayaxioms{} axioms are all declared in Lean's core library; the archive declares none,
and the five theorems depend on three of them\ifsecondrun\else, by the recorded lists of
\cref{sec:axioms}\fi. In \replayrepeated{} cases a name is declared
again by a further module, which Lean's import accepts for theorems with the same statement;
every version was checked. The \replayskipped{} declarations not replayed are those that Lean
marks unsafe or partial: \replayskippedrec{} are the compiled forms of recursive definitions
and the rest are programs and the unsafe axioms that Lean's compiler uses. Lean's own replay
does not check them either, and outside Lean's
module system the kernel accepts no theorem or safe definition that refers to one;
\cref{sec:checker} says what holds inside it.}

\pagestyle{empty}
\setlength{\parindent}{0pt}
\setlength{\parskip}{0pt}
\linespread{0.96}
\begin{document}
\small
\begin{center}
{\large\bfseries The argument of Alp\"oge's formal proof of the Riemann Hypothesis, on one page}\\[2pt]
{\footnotesize Qui Vincit and Fable. \paperdate, with version \paperversion{} of the paper
\emph{The argument of Alp\"oge's formal proof of the Riemann Hypothesis: an account read from the
Lean source}, \url{https://github.com/vincitamore/riemann-lean-account}.}
\end{center}
\vspace{-2pt}
{\footnotesize Each step is a statement of the Lean archive, named in brackets and read in the
source. The paper's Section 2 says what the machine checked; this page says what the whole
is. Letters are those of the paper.}

\step{Theorem} [\lean{LFunction_ne_zero_of_half_lt_re}, \lean{riemannHypothesis}]\textbf{.}
\emph{For every Dirichlet character $\chi$ and every $s$ with $\re s>\tfrac12$, other than the
pole of $\zeta$ at $1$, $L(s,\chi)\ne0$. Hence every zero of $\zeta$ other than the trivial ones
has real part $\tfrac12$.}

\step{Proof.} It suffices to treat primitive characters; the second sentence follows by the
functional equation. Fix a small $\eps>0$. Below, $x$ is a large scale and $b=\exp(x^\eps)$.

\step{1. Landau on windows} [\lean{zeroFree_of_rpow_window}]\textbf{.} For a character $\chi$ put
$G_\chi(y)=\sum_{n\le y}\Lambda(n)\chi(n)(\tfrac yn-1)+y\tfrac{L'}{L}(1,\chi)$, with
$-y\log y+(1+\gamma)y$ in place of the last term when $\chi=1$. Let $\mathcal K$ be a finite set
of primitive characters, closed under inversion and containing $1$, and
$G_{\mathcal L}=G_1+\sum_{\psi\in\mathcal K,\,\psi\ne1}\mathrm{wt}(\psi)\re G_\psi$ with positive
weights. If for arbitrarily large $Y$ one has $G_{\mathcal L}(y)\le y^{1/2+C\eps}$ for all
$y\in[Y,Y^{3/2}]$, then no $L(s,\psi)$, $\psi\in\mathcal K$, vanishes for
$\re s>\tfrac12+C\eps$. So it is enough to bound $G_{\mathcal L}(x)$ from above at the scales of
such windows.

\step{2. An exact identity} [\lean{ch_balance}]\textbf{.} Let $E$ be an abelian CM field (a
\emph{host}) with Galois group $\mathfrak G=(\Z/2)^2\times H$ and CM type $\Phi$, a subgroup of
index two, ramified only at primes whose inertia groups and Frobenius elements lie in $\Phi$
and at one prime $d$, whose inertia group has order two and is not in $\Phi$. For a size $a>0$,
weights $w>0$ on $\Phi$, and a pairing $g\mapsto\rho_g$
of $\mathfrak G\setminus\Phi$ with $\Phi$, put
\[
\ASK=\sum_p2\log p\!\!\sum_{i\ge1,\ \Frob_p^i\notin\Phi}\!\!\max\!\bigl(\tfrac a{p^i}-w(\rho_{p,i}),0\bigr)+(\text{a term at }d),
\qquad
\MAIN=a\Bigl(\log a-1-\gamma+\tfrac{L'}{L}(1,\chi_\Phi)-\tfrac1{\abs\Phi}\!\sum_{\rho\in\Phi}\!\log w(\rho)\Bigr).
\]
Then $\ASK-\MAIN=\frac1{\abs\Phi}\sum_{g\notin\Phi}\sum_{\psi}\overline\psi(g)\,w(\rho_g)\,G_\psi\bigl(a/w(\rho_g)\bigr)$,
the inner sum over the characters of $\mathfrak G$. Nothing is estimated.

\step{3. From the count to $G$} [\lean{ch_core}, \lean{core_gm}, \lean{GrowthSizeBound}]\textbf{.}
Take sizes $a_t$ with $\abs{a_t/b-x}\le1$ for $t$ in a set $TT\subset H$, and $m$ systems of
weights $w_{u_0,j}$ for each translate $u_0$: about $b$ on a thick set of density
$\delta_T\ge x^{-\eps}$, so that $a/w\approx x$, and so small elsewhere that $a/w$ is a scale at
which the classical zero-free region makes every $G_\psi$ negligible. If for every $u_0$
\[
\sum_{t\in TT}\sum_{j\le m}(\ASK-\MAIN)(a_t,t,w_{u_0,j})\ \le\ x^{C\eps}\,\Br,\qquad
\Br=\abs{TT}\,a_E\bigl(1+\tfrac{m^2}k\bigr)+m\bigl(2nmb\,x^\eps+\abs{T_2}\,k\,b\,x^{2\eps}/n^2\bigr),\quad a_E\le2xb,
\]
then, the sizes $TT$ being equidistributed for the kept characters,
$G_{\mathcal L}(x)\le x^{1/2+(C+3)\eps}+\sum\abs{G_\psi(x)}$, the sum over the characters of
the host that are not kept. Divided by the normalisation $\abs{TT}\,m\,b\,\delta_T$ the three
terms of $\Br$ are of sizes $x/m$, $m$ and $k/n^2$ up to $x^{O(\eps)}$. With
$m\approx\sqrt x$ lifts and $xm\le k\le n^2\sqrt x$ all three are $x^{1/2+O(\eps)}$.
\emph{This is the only place where $\tfrac12$ is produced.}

\step{4. The upper bound} [\lean{each_unit_bound_statement}; the record \lean{SizeData}\textsubscript{\ttfamily\footnotesize 2}]\textbf{.}
For each of many words $u$ there are: a Hermitian lattice $M_u$ of algebraic modular forms on a
definite unitary group over $E$; the rescaling
$N_u=\prod_\lambda\fP_\lambda^{-\lfloor e_\lambda g(u,\lambda)\rfloor}\cdot(\text{a sublattice of }M_u)$,
the exponent $g(u,\lambda)$ a sum over the $m$ lifts of a \emph{gain} and two corrections; and a
map $B_u$ of height at most $h(u)$, with $B_u(N_u)\subseteq P$ for one lattice $P$
[\lean{adm2}], sending a form to the periods of its $m$-fold iterated theta lift. Four clauses
make the chain, $\mu$ denoting the slope:
\[
\begin{aligned}
&\text{(Loss)}\ \ \ASK-\mathrm{Ga}\le\mathrm{Lo};
&&\text{(Gain)}\ \ \mu(M_u)+\mathrm{Ga}+\mathrm{Fo}-\mathrm{Ee}(u)-(\text{loss})\le\mu(N_u);\\
&\text{(Slope)}\ \ \overline{\mu(N_u)-h(u)}-(\text{losses})\le\mu(P);
&&\text{(Div)}\ \ \mu(P)+h(u)-\mathrm{Fo}\le\textstyle\sum_{t}\bigl(m(\ph(a_t)-a_t\log\Omega)-a_t\log T(u)\bigr)+(\text{small}).
\end{aligned}
\]
Here $\mathrm{Ga}$ is the total gain net of a cost, $\mathrm{Fo}$ and $\mathrm{Ee}(u)$ are the
totals of the corrections, the bar is the rank-weighted mean over $u$,
$\ph(y)=y\log y-y$, $\Omega$ is a CM period and $T(u)$ a level prime. A dimension count bounds
$\log T(u)$ from below [\lean{Count5}]; then the right side of (Div) is at most
$\MAIN-m\ell_g\sum_ta_t$ up to an error of the first kind, by an identity
[\lean{key_identity5}], and $\mathrm{Ee}(u)$ is $m\ell_g\sum_ta_t$ up to an error
[\lean{EComp5}]; $\ell_g$ is a constant of the host. Chaining gives the hypothesis of step 3.

\step{5. Why the four clauses hold.}
\emph{Loss and Gain: divisibility at inert primes} [\lean{gainP}, \lean{adm2},
\lean{admissible_level_two}, \lean{CM2}, \lean{CM3}, \lean{IntInert}, \lean{LemmaSt},
\lean{LemmaRA}]. At a prime $\ell\ge5$ inert in
$\Q(\sqrt{-d})$ a CM elliptic curve has good reduction and its formal group is the Lubin--Tate
group of height two with logarithm $\sum_jX^{\ell^{2j}}/\ell^j$. For a series integral in the
parameter, $r!$ times the $r$-th Taylor coefficient in the uniformising coordinate has
$\ell$-adic valuation at least $\gamma_\ell r-(\log_\ell r+1)$,
$\gamma_\ell=\ell/(\ell^2-1)=\sum_{e\ \mathrm{odd}}\ell^{-e}$. The kernel of the theta lift is
such a series, so the periods are divisible to the order $\lfloor e_\lambda g(u,\lambda)\rfloor$,
and the gain is a closed formula whose main term is the odd-power count of step 2.
\emph{Slope} [\lean{TheoremF}, \lean{InputSG}, \lean{KExpansion}, \lean{UniqGap},
\lean{SiegelDensities}, \lean{pbApplied}]. A slope inequality for nearly tight frames of maps
between Hermitian lattices bounds $\mu(P)$ from below by the mean of $\mu(N_u)-h(u)$. The maps
$B_u$ are compressions of Bessel--Poincar\'e series, and a relative trace formula makes their
Gram blocks nearly those of independent random subspaces.
\emph{Div} [\lean{BoundsMu}, \lean{divisor_inequality_assembly}, \lean{lcv_le},
\lean{topCovolSlim2_use_alt_alt}, \lean{chowla_selberg_eta}]. Determinants factor along the $m$
lifts; a ledger over the lifts cancels every factorial and every period but the top ones; the
covolumes of the auxiliary lattices of integral forms are bounded by arithmetic intersection
theory on a Shimura variety; and the mean of $\log\Omega$ is given by the Chowla--Selberg
formula, which is how $\tfrac{L'}{L}(1,\chi_\Phi)$ reappears.

\step{6. The inputs} [\lean{every_character_of_green_v_win}, \lean{supplyChi}, \lean{PairEx}]\textbf{.}
Hosts with these properties exist at every scale of long windows, the characters that are not
kept being bounded at the single point $x$ by a choice of the two auxiliary primes. This needs:
(i)~a bound on Galois cohomology [\lean{VWin}]:
$\sum_\ell\log\ell\sum_{1\le e<\lfloor x^{401}\rfloor}\len\Sha(E,\ell,e)\le\exp(x^\eps)^{1/2}$ on windows,
proved through congruences between cusp forms and Eisenstein classes on $\U(3)$;
(ii)~Green's parametrisation of the cuspidal characters of $\GL_\nu(\F_p)$ [\lean{Green}];
(iii)~\rowspublished{} published theorems, each proved in the archive (the paper, Appendix~C). \hfill$\square$

\step{What this page is.} A statement of the skeleton, not a proof that can be checked by hand.
Every bracketed name below the theorem is a declaration that the final theorem depends on. The page
omits the definitions of the lattices, of the gain and of the error terms, the size conditions,
and the proofs; the paper gives the first three and says where each of the rest was read.
\end{document}
